\documentclass[a4paper,12pt]{article}

\usepackage[T1]{fontenc}
\usepackage[italian]{babel}
\usepackage{graphicx}
\usepackage{geometry}
\usepackage{tabularx}
\usepackage[table]{xcolor}
\usepackage[most]{tcolorbox}
\usepackage{fancyhdr}
\usepackage[hidelinks]{hyperref}

\newcommand{\groupmail}{alittlebyte.swe@gmail.com}
\newcommand{\grouplogo}{Immagini/Logo_Progetto-ezgif.com-crop.png}
\newcommand{\groupsymbol}{Immagini/solo_logo.png}

\geometry{
    left=2.5cm,
    right=2.5cm,
    top=2.5cm,
    bottom=2.5cm,
    headheight=331pt,
    includeheadfoot
}

\pagestyle{fancy}
\fancyhf{}

\renewcommand{\headrulewidth}{0pt}
\fancyhead{}

\renewcommand{\footrulewidth}{0.4pt}
\fancyfoot[L]{\includegraphics[height=0.6cm]{\groupsymbol}\hspace{0.45em}aLittleByte}
    \fancyfoot[C]{\thepage}
\fancyfoot[R]{Verbale 30/03/2026}

\fancypagestyle{firstpage}{
    \fancyhf{}
    \renewcommand{\headrulewidth}{0pt}
    \renewcommand{\footrulewidth}{0.4pt}
    \fancyhead[C]{%
        \makebox[\textwidth][c]{%
            \begin{tabular}{c}
                \includegraphics[height=10cm]{\grouplogo}\\[1ex]
                \small\texttt{\groupmail}
            \end{tabular}%
        }
    }
    \fancyfoot[L]{\includegraphics[height=0.6cm]{\groupsymbol}\hspace{0.45em}aLittleByte}
        \fancyfoot[C]{\thepage}
    \fancyfoot[R]{Verbale 30/03/2026}
}

\begin{document}

\thispagestyle{firstpage}

\vspace*{0.5cm}

\begin{center}
    {\LARGE \textbf{Verbale Interno di Riunione - 30/03/2026}}
\end{center}

\vspace{1.5cm}

\begin{center}
    \renewcommand{\arraystretch}{1.3}
    \begin{tcolorbox}[
        width=0.92\textwidth,
        colback=gray!6,
        colframe=gray!45,
        boxrule=0.5pt,
        arc=2mm,
        boxsep=0pt,
        left=8pt, right=8pt, top=8pt, bottom=8pt
    ]
        \begin{tabularx}{\linewidth}{>{\bfseries}p{3.5cm} X}
            Gruppo      & aLittleByte (gruppo 19) \\
            Redattore   & Alex Oloni \\
            Verificatori & V.Y. Kaneda Fernandes\\
            Destinatari & Prof. Tullio Vardanega, Prof. Riccardo Cardin \\
            Data        & 30 Marzo 2026\\
        \end{tabularx}
    \end{tcolorbox}
\end{center}

\newpage

\newgeometry{left=2.5cm,right=2.5cm,top=2.5cm,bottom=2.5cm,includefoot}
\tableofcontents
\newpage
\section{Informazioni Generali}
\section*{Specifiche Documento}
\hrule
\vspace{1cm}

\begin{tabular}{>{\bfseries}p{5cm} | p{10cm}}
Luogo Riunione & Discord \\[6pt]

Orario (Inizio - Fine) & 10:00 -- 10:45 \\[6pt]

Redattore &
  A. Oloni\\[6pt]

Amministratore & 
A. Gastaldello\\[6pt]

Verificatore &
V. Y. Kaneda Fernandes\\[6pt]

Partecipanti &
  A. Oloni\\ &
  L. Soligo\\&
  A. Gastaldello\\&
  E. Bellet\\&
  D. Belluz\\&
  A. Maule\\& 
  V. Y. Kaneda Fernandes\\[6pt]


\end{tabular}


\section{Ordine del Giorno}
    \begin{itemize}
    \item Pianificazione contatto iniziale con l'azienda
    \item Assegnazione e organizzazione del lavoro sui documenti di progetto
    \end{itemize}


\section{Attività svolte}
\subsection{Pianificazione contatto iniziale con l'azienda}
Il gruppo ha discusso se fosse meglio contattare l'azienda tramite mail o colloquio. 
Si è deciso preliminarmente di inviare una mail per presentarsi e chiedere consigli su come procedere. L'invio della mail è
previsto entro lo stesso pomeriggio della giornata della riunione. Si è stabilito che sarà Angelica Gastaldello a occuparsi del contatto e della comunicazione con l'azienda, 
affinché venga informata sulla modalità di lavoro adottata dal gruppo e si possano ricevere suggerimenti o approvazioni.
\subsection{Assegnazione e organizzazione del lavoro sui documenti di progetto}
È stata fatta una revisione e divisione dei documenti da preparare: 
analisi dei requisiti, norme di progetto, glossario, lista dei quesiti e mock-up. Vista la complessità e la lunghezza di alcuni documenti, 
si è convenuto che l'analisi dei requisiti sarà sviluppata da due persone (Angelica e Angela) per evitare perdita di tempo e facilitarne la stesura. 
Le norme di progetto saranno gestite da Diego ed Eleonora, mentre Yuri si occuperà del glossario. 
È stato discusso che il piano di progetto e il piano di qualifica saranno affrontati più avanti, considerati meno immediati. 
Per i mock-up si è deciso di iniziare presto, coinvolgendo Leonardo Soligo e Alex Oloni per ottenere prima un prodotto da sottoporre all'azienda per una validazione preliminare.
Quindi i documenti che saranno iniziati a scrivere saranno:
\begin{itemize}
    \item Analisi dei requisiti
    \item Norme di progetto
    \item Glossario
    \item Piano di Progetto
\end{itemize}

\subsection{Scelta degli strumenti per la realizzazione dei mock-up}
Si è discusso l'uso di strumenti per la progettazione dei mock-up, con un confronto tra Canva e Figma. 
Figma è stato indicato come preferibile per la sua versatilità, disponibilità di librerie e possibilità di collaborazione, 
anche se nella versione gratuita con alcune limitazioni. 
È stato menzionato un'alternativa open source, Penpot, che gestisce file SVG, permettendo la gestione delle modifiche tramite commit in repository. 
Si è concordato di utilizzare Penpot come programma per il mock-up.

\section{To-Do List}
\begin{table}[h]
    \centering
    \begin{tabular}{|p{4.5cm}|p{5cm}|l|}
    \hline
        \rowcolor{lightgray}\textit{\textbf{Persona Incaricata}}  & \textbf{Task Assegnata} &\textbf{Link Issue} \\
    \hline
        \textit {Alex Oloni e Leonardo Soligo} & Creazione Mock-up & \href{https://github.com/aLittleByte-19/Documentazione/issues/50}{\#50}\\
    \hline
         \textit {Angelica Gastaldello e Angela Maule} & Inizio stesura Analisi dei Requisiti& \href{https://github.com/aLittleByte-19/Documentazione/issues/51}{\#51}\\
    \hline
        \textit{Diego Belluz e Eleonora Bellet} &  Inizio stesura Norme di Progetto& \href{https://github.com/aLittleByte-19/Documentazione/issues/52}{\#52}\\
    \hline
        \textit{Vinicius Yuri Kaneda Fernandes} & Inizio stesura Glossario & \href{https://github.com/aLittleByte-19/Documentazione/issues/53}{\#53}\\
    \hline
        \textit{Alex Oloni} & Stesura Verbale Interno 31/03/2026 & \href{https://github.com/aLittleByte-19/Documentazione/issues/54}{\#54}\\
    \hline
        \textit {Angelica Gastaldello}  &  Scrittura email all'azienda & \href{https://github.com/aLittleByte-19/Documentazione/issues/55}{\#55}\\
    \hline
    \end{tabular}
    \label{tab:placeholder}
\end{table}

\end{document}
